\chapter{Μεθοδολογία - Μοντέλο Προσομοίωσης}
\label{ch:methodology}

\section{Επισκόπηση}
\label{sec:meth-overview}

Το σενάριο αναφοράς αποτελείται από δύο επίγειους σταθμούς βάσης (BS1,
BS2), έναν δορυφόρο LEO (SAT-1) σε υψόμετρο $550$~km, και έξι χρήστες,
τοποθετημένους σε πραγματικές γεωγραφικές συντεταγμένες γύρω από την
περιοχή της Αθήνας (τρεις χρήστες κοντά στους σταθμούς βάσης, τρεις σε
απόσταση 100+~km, όπου μόνο ο δορυφόρος προσφέρει ρεαλιστική κάλυψη). Για
κάθε χρήστη, η προσομοίωση υπολογίζει το SNR κάθε υποψήφιας ζεύξης (προς
κάθε σταθμό βάσης, και προς τον δορυφόρο), επιλέγει ως εξυπηρετητή τον
κόμβο με το υψηλότερο SNR (single-connectivity), και εξάγει χωρητικότητα
και ενέργεια ανά bit για τον νικητή. Συνοπτικά, η ροή υπολογισμού είναι:
γεωμετρία $\rightarrow$ path loss (ανά υποψήφιο) $\rightarrow$ SNR (ανά
υποψήφιο) $\rightarrow$ επιλογή κόμβου $\rightarrow$ κατανομή εύρους
ζώνης $\rightarrow$ χωρητικότητα/ενέργεια.

Ο πυρήνας αυτής της λογικής υλοποιείται στη συνάρτηση \texttt{simulateScenario.m}
(βλ.~Παράρτημα~Ι). Η συνάρτηση αυτή καλείται είτε μία φορά (στατικό
σενάριο αναφοράς, \texttt{test\_simulation.m}), είτε επαναληπτικά πάνω σε
τυχαιοποιημένες τοπολογίες (\texttt{monteCarloDriver.m}), είτε
επαναληπτικά με την ίδια τοπολογία αλλά διαφορετική στοχαστική
πραγματοποίηση καναλιού (\texttt{kpiRepeatedRuns.m}), είτε επαναληπτικά
με κινούμενο δορυφόρο στον χρόνο (\texttt{temporalPassSimulation.m}). Οι
τρεις τελευταίοι driver scripts μοιράζονται τον ίδιο πυρήνα υπολογισμού
link-budget· διαφοροποιούνται μόνο ως προς το τι μεταβάλλουν μεταξύ
διαδοχικών κλήσεων (γεωμετρία, seed τυχαιότητας, ή/και θέση δορυφόρου).

\section{Μοντέλο επίγειου καναλιού}
\label{sec:meth-terrestrial}

Για κάθε ζεύξη χρήστη-σταθμού βάσης, υπολογίζεται πρώτα η 3D απόσταση
$d_{3D}$ (μέσω μετατροπής γεωδαιτικών συντεταγμένων σε τοπικό ENU πλαίσιο
αναφοράς, WGS84 ελλειψοειδές) και η οριζόντια απόσταση $d_{2D}$.

\subsection{Στοχαστική επιλογή LOS/NLOS}

Αντί για μία σταθερή, καθολική παραδοχή οπτικής επαφής, κάθε ζεύξη
αποφασίζεται ανεξάρτητα ως LOS ή NLOS με πιθανότητα $p_{\text{LOS}}(d_{2D})$
όπως ορίζεται στο TR~38.901 §7.4.2 (Πίνακας 7.4.2-1). Για το σενάριο UMa,
που χρησιμοποιείται στην παρούσα εργασία:
\begin{equation}
p_{\text{LOS}}(d_{2D}) =
\begin{cases}
1, & d_{2D} \le 18\text{~m} \\[4pt]
\left(\dfrac{18}{d_{2D}} + e^{-d_{2D}/63}\left(1 - \dfrac{18}{d_{2D}}\right)\right)
\left(1 + C'(d_{2D}, h_{\text{UT}})\right), & d_{2D} > 18\text{~m}
\end{cases}
\label{eq:los-prob-uma}
\end{equation}
όπου $h_{\text{UT}}$ το ύψος του χρήστη και $C'$ ένας διορθωτικός όρος που
μηδενίζεται για $h_{\text{UT}} \le 13$~m (η συνήθης περίπτωση πεζού/κινητού
χρήστη). Η απόφαση LOS/NLOS προκύπτει τότε από μια δοκιμή Bernoulli:
$\text{LOS} \Leftrightarrow U < p_{\text{LOS}}(d_{2D})$, με $U \sim
\mathcal{U}(0,1)$ ανεξάρτητο ανά ζεύξη και ανά εκτέλεση.

\subsection{Path loss και shadow fading}

Δεδομένης της κατάστασης LOS/NLOS, η μέση απώλεια διάδοσης
$\overline{PL}$ υπολογίζεται μέσω της συνάρτησης \texttt{nrPathLoss} του
5G Toolbox (μοντέλο TR~38.901, σενάριο UMa ή UMi), η οποία επιστρέφει
επίσης την τυπική απόκλιση shadow fading $\sigma_{SF}$ του σεναρίου/κατάστασης.
Η τελική (στοχαστική) απώλεια διάδοσης είναι
\begin{equation}
PL = \overline{PL} + \sigma_{SF} \cdot Z, \qquad Z \sim \mathcal{N}(0,1),
\label{eq:pathloss-shadowing}
\end{equation}
όπου, στο στατικό σενάριο αναφοράς, κάθε ζεύξη σε κάθε εκτέλεση της
προσομοίωσης αντλεί το δικό της, ανεξάρτητο δείγμα shadow fading (στη
χρονική προσομοίωση αυτό τροποποιείται· βλ.
Ενότητα~\ref{subsec:meth-correlated-sf}).

\subsection{SNR}

Το SNR της ζεύξης προς σταθμό βάσης υπολογίζεται σε dB ως
\begin{equation}
\text{SNR}_{\text{BS}}[\text{dB}] = (P_{Tx} - 30) - PL - N_{\text{BS}}[\text{dBW}],
\label{eq:snr-bs}
\end{equation}
όπου $P_{Tx} = 43$~dBm η ισχύς εκπομπής ανά σταθμό βάσης, και
$N_{\text{BS}}$ η ισχύς θορύβου θερμικής προέλευσης,
\begin{equation}
N_{\text{BS}}[\text{dBW}] = 10\log_{10}\!\big(k_B \, T_{\text{eq}} \, B_{\text{BS}}\big),
\qquad T_{\text{eq}} = T_0 + 290(F - 1),
\label{eq:noise-power}
\end{equation}
με $k_B$ τη σταθερά Boltzmann, $T_0 = 290$~K τη θερμοκρασία κεραίας λήψης,
$F$ τον γραμμικό λόγο του noise figure δέκτη ($5$~dB), και $B_{\text{BS}}$
το εύρος ζώνης του φέροντος (υπολογισμένο από τον αριθμό Resource Blocks
του NR carrier configuration: $51$ RB $\times 12$ subcarriers $\times
30$~kHz subcarrier spacing $\approx 18.4$~MHz).

\section{Μοντέλο δορυφορικού καναλιού}
\label{sec:meth-sat-channel}

Για τη ζεύξη χρήστη-δορυφόρου, υπολογίζεται πρώτα η γωνία ανύψωσης
(elevation) και η slant range απόσταση μέσω μετατροπής γεωδαιτικών
συντεταγμένων σε Azimuth-Elevation-Range (AER). Η ζεύξη θεωρείται ορατή
μόνο εάν η γωνία ανύψωσης υπερβαίνει ένα ελάχιστο όριο ορατότητας
$\theta_{\min} = 10^\circ$, μια συνήθη τιμή μάσκας ορατότητας δορυφόρου
στη βιβλιογραφία NTN. Όταν ορατή, η απώλεια διάδοσης υπολογίζεται ως
απώλεια ελεύθερου χώρου (FSPL):
\begin{equation}
PL_{\text{sat}}[\text{dB}] = 20\log_{10}\!\left(\frac{4\pi \, d_{\text{slant}}}{\lambda}\right),
\qquad \lambda = \frac{c}{f_c},
\label{eq:fspl}
\end{equation}
με $f_c = 2.01$~GHz (ζώνη S) και $d_{\text{slant}}$ τη slant range
απόσταση. Το SNR της δορυφορικής ζεύξης είναι τότε
\begin{equation}
\text{SNR}_{\text{sat}}[\text{dB}] = (\text{EIRP} - 30) - PL_{\text{sat}} - N_{\text{sat}}[\text{dBW}],
\label{eq:snr-sat}
\end{equation}
όπου $\text{EIRP} = P_{Tx,\text{sat}} + G_{\text{ant}} = 34 + 30 =
64$~dBm (ισχύς ενισχυτή $34$~dBm, κέρδος κατευθυντικής κεραίας $30$~dBi,
σύμφωνα με τυπικές τιμές TR~38.821 για δορυφόρο LEO), και $N_{\text{sat}}$
η ισχύς θορύβου του δέκτη υπολογισμένη όπως στην~\eqref{eq:noise-power}
αλλά με το εύρος ζώνης του δορυφόρου ($20$~MHz). Όταν η γωνία ανύψωσης
είναι μικρότερη του ορίου ορατότητας, η ζεύξη θεωρείται μη-διαθέσιμη
($PL_{\text{sat}} = \infty$, $\text{SNR}_{\text{sat}} = -\infty$).

Πρόκειται για ένα ηθελημένα απλοποιημένο μοντέλο: δεν λαμβάνει υπόψη
ατμοσφαιρικές ή σπινθηριστικές απώλειες (TR~38.821 §6.1), ούτε κέρδος
κεραίας δέκτη εξαρτώμενο από τη γωνία. Πιο σημαντικό για τα ευρήματα του
Κεφαλαίου~\ref{ch:results} είναι ότι δεν εμπεριέχει καμία στοχαστική
συνιστώσα (fading) στη δορυφορική ζεύξη· η δορυφορική ζεύξη είναι, στην
τρέχουσα προσομοίωση, πολύ πιο προβλέψιμη από την επίγεια, κάτι που
πρέπει να ληφθεί υπόψη κατά την ερμηνεία των αποτελεσμάτων του
Κεφαλαίου~\ref{ch:results}, και αναφέρεται ως κατεύθυνση μελλοντικής
εργασίας στο Κεφάλαιο~\ref{ch:discussion}.

\section{Επιλογή κόμβου εξυπηρέτησης}
\label{sec:meth-node-selection}

Για κάθε χρήστη $u$, ο κόμβος εξυπηρέτησης επιλέγεται ως
\begin{equation}
n^\star(u) = \argmax_{n \in \{\text{BS}_1, \ldots, \text{BS}_{N_{\text{BS}}}, \text{SAT-1}\}} \text{SNR}_n(u),
\label{eq:node-selection}
\end{equation}
δηλαδή μια \emph{greedy}, ανά χρήστη επιλογή του υποψήφιου με το υψηλότερο
στιγμιαίο SNR, χωρίς καμία μορφή υστέρησης (hysteresis) ή περιθωρίου
απόφασης. Είναι single-connectivity σχήμα: κάθε χρήστης συνδέεται σε
έναν και μόνο κόμβο. Η επέκταση προς πραγματική πολυσυνδεσιμότητα (π.χ.
ταυτόχρονη εξυπηρέτηση από BS \emph{και} δορυφόρο, όπως στο Mode~1 της
αρχιτεκτονικής SS-SBS των~\cite{li2025advancing}), σε αντίθεση με τα
σχήματα που συζητούνται στο Κεφάλαιο~\ref{ch:background}, παραμένει
ανοιχτή κατεύθυνση μελλοντικής εργασίας.

Ο κανόνας~\eqref{eq:node-selection} δεν προβλέπει κατάσταση
\emph{outage}: ακόμα και όταν και οι δύο υποψήφιοι
προσφέρουν αμελητέο ή αρνητικό SNR (π.χ. χρήστης εκτός εμβέλειας κάθε
σταθμού βάσης, με τον δορυφόρο ταυτόχρονα κάτω από το όριο ορατότητας), ο
χρήστης θα ανατεθεί στον "λιγότερο κακό" υποψήφιο. Το φαινόμενο αυτό
τεκμηριώνεται εμπειρικά στην Ενότητα~\ref{sec:results-temporal}.

\section{Κατανομή εύρους ζώνης και χωρητικότητα}
\label{sec:meth-capacity}

Το συνολικό εύρος ζώνης κάθε κόμβου ($B_{\text{BS}}$ για σταθμό βάσης,
$B_{\text{sat}} = 20$~MHz για τον δορυφόρο) κατανέμεται ισομερώς μεταξύ
όλων των χρηστών που εξυπηρετεί ο συγκεκριμένος κόμβος (\emph{node load}
$L_n$):
\begin{equation}
B_u = \frac{B_n}{L_n}, \qquad L_n = \big|\{u' : n^\star(u') = n\}\big|.
\label{eq:bw-split}
\end{equation}
Η χωρητικότητα ανά χρήστη υπολογίζεται τότε μέσω του τύπου Shannon πάνω
στο μερίδιο εύρους ζώνης του χρήστη και το SNR του νικητή κόμβου:
\begin{equation}
C_u = B_u \log_2\!\big(1 + \text{SNR}_{n^\star(u)}(u)\big) \quad [\text{bit/s}].
\label{eq:shannon-capacity}
\end{equation}
Πρόκειται και εδώ για μια αισιόδοξη, μη-κβαντισμένη (unbounded Shannon)
προσέγγιση της χωρητικότητας. Ο περιορισμός της σε πραγματικά επίπεδα
MCS/CQI (TS~38.214~\cite{ts38214}) αναφέρεται ως κατεύθυνση μελλοντικής
εργασίας.

\section{Ενεργειακό μοντέλο και ενέργεια ανά bit}
\label{sec:meth-energy}

Η κατανάλωση ισχύος κάθε κόμβου υπολογίζεται με διαφορετικό μοντέλο ανά
τύπο κόμβου. Για τον επίγειο σταθμό βάσης, υιοθετείται το γραμμικό μοντέλο
EARTH~\cite{auer2011earth}:
\begin{equation}
P_{\text{BS}} = N_{\text{TRX}}\big(P_0 + \Delta_P \, P_{\text{out}}\big),
\label{eq:earth-power}
\end{equation}
με $N_{\text{TRX}} = 1$ (μονο-sector μοντέλο), $P_0 = 130$~W (σταθερή
κατανάλωση σε ενεργή λειτουργία), $\Delta_P = 4.7$ (κλίση κατανάλωσης ως
προς την εκπεμπόμενη ισχύ), και $P_{\text{out}}$ η γραμμική εκπεμπόμενη
ισχύς ($43$~dBm $= 20$~W). Για τον δορυφόρο, υιοθετείται ένα απλό γραμμικό
μοντέλο απόδοσης ενισχυτή ισχύος (PA):
\begin{equation}
P_{\text{sat}} = P_{\text{fix}} + \frac{P_{\text{out,sat}}}{\eta_{\text{PA}}},
\label{eq:pa-power}
\end{equation}
με $P_{\text{fix}} = 0$~W (συντηρητική παραδοχή -- δεν μοντελοποιείται
ξεχωριστά η κατανάλωση του υπόλοιπου payload) και $\eta_{\text{PA}} =
0.4$ (τυπική τιμή απόδοσης SSPA/TWTA, εντός του εύρους $0.35$-$0.5$).

Η ισχύς κάθε κόμβου κατανέμεται επίσης ισομερώς στους $L_n$ χρήστες που
εξυπηρετεί (ίδια λογική με το εύρος ζώνης), και η ενέργεια ανά bit
προκύπτει διαιρώντας με τον ρυθμό bit του χρήστη:
\begin{equation}
E_u \, [\mu\text{J/bit}] = \frac{P_n / L_n}{C_u} \times 10^6.
\label{eq:energy-per-bit}
\end{equation}
Το μοντέλο αυτό είναι proxy, όχι μέτρηση υλικού: δεν λαμβάνει υπόψη την
κατανάλωση κόμβων σε αδράνεια (χωρίς εξυπηρετούμενους χρήστες), παρά
μόνο ενεργών κόμβων.

\section{Μοντέλο κίνησης δορυφόρου LEO για τη χρονική προσομοίωση}
\label{sec:meth-temporal}

Το στατικό σενάριο αναφοράς (Ενότητες~\ref{sec:meth-terrestrial}-\ref{sec:meth-energy})
δεν περιλαμβάνει καμία χρονική εξέλιξη: η θέση του δορυφόρου, όπως και οι
θέσεις σταθμών βάσης/χρηστών, παραμένουν σταθερές. Ο επιβλέπων καθηγητής
ζήτησε να εξεταστεί πώς συμπεριφέρεται το μοντέλο επιλογής κόμβου καθώς ο
δορυφόρος LEO κινείται πάνω από την περιοχή κάλυψης, ώστε να αποτυπωθεί
τόσο η κίνησή του όσο και η μεταβολή του καναλιού σε διαδοχικές
μεταδόσεις. Για τον σκοπό αυτό αναπτύχθηκε ένα απλοποιημένο μοντέλο
ίχνους εδάφους (ground track).

\subsection{Ταχύτητα ίχνους εδάφους}

Η ταχύτητα με την οποία το υποδορυφορικό σημείο διασχίζει την επιφάνεια
της Γης υπολογίζεται από την Κεπλεριανή τροχιακή περίοδο μιας κυκλικής
τροχιάς σε υψόμετρο $h$:
\begin{equation}
T_{\text{orbit}} = 2\pi\sqrt{\frac{a^3}{\mu}}, \qquad a = R_\oplus + h,
\label{eq:orbital-period}
\end{equation}
με $\mu = 3.986004418\times10^{14}$~m$^3$/s$^2$ τη βαρυτική παράμετρο της
Γης και $R_\oplus = 6371$~km τη μέση ακτίνα της. Η ταχύτητα ίχνους εδάφους
προκύπτει τότε ως
\begin{equation}
v_{\text{ground}} = \frac{2\pi}{T_{\text{orbit}}} \, R_\oplus,
\label{eq:ground-speed}
\end{equation}
μια προσέγγιση που αγνοεί τη συνεισφορά της ιδίας περιστροφής της Γης
(δηλαδή υπολογίζει την ταχύτητα του ίχνους ως προς μη-περιστρεφόμενη Γη).
Για $h = 550$~km, οι εξισώσεις~\eqref{eq:orbital-period}-\eqref{eq:ground-speed}
δίνουν $T_{\text{orbit}} \approx 5730$~s ($\approx 95.5$~λεπτά) και
$v_{\text{ground}} \approx 6986$~m/s, τιμές συνεπείς με πραγματικούς
δορυφόρους LEO στο συγκεκριμένο υψόμετρο.

\subsection{Ίχνος εδάφους και χρονικός βρόχος προσομοίωσης}

Το υποδορυφορικό σημείο μετακινείται σε σταθερό γεωγραφικό πλάτος (ίσο με
το πλάτος του κέντρου του BS cluster) προς ανατολάς, με ταχύτητα
$v_{\text{ground}}$:
\begin{equation}
\text{lon}_{\text{sub}}(t) = \text{lon}_{\text{center}} +
\frac{\Delta x_0 + v_{\text{ground}} \, t}{111.32 \cdot \cos(\text{lat}_{\text{center}}) \; [\text{km/}^\circ]},
\label{eq:ground-track}
\end{equation}
όπου $\Delta x_0$ η αρχική μετατόπιση (km) ώστε ο δορυφόρος να ξεκινά
αόρατος/οριακά ορατός. Πρόκειται για μια απλοποίηση επίπεδης γεωμετρίας
(όχι πλήρη ορβιτογράφο SGP4 ή γεωδαιτικό υπολογισμό great-circle), αρκετή
ωστόσο ώστε το προκύπτον προφίλ γωνίας ανύψωσης να αναπαράγει πιστά το
αναμενόμενο σχήμα μίας διέλευσης δορυφόρου (άνοδος από το όριο ορατότητας
έως σχεδόν το ζενίθ και επάνοδος), όπως επιβεβαιώνεται εμπειρικά στην
Ενότητα~\ref{sec:results-temporal}.

Σε κάθε χρονικό βήμα $t_k = k \cdot \Delta t$ (η προσομοίωση χρησιμοποιεί
$\Delta t = 5$~s), υπολογίζεται η νέα θέση δορυφόρου μέσω
της~\eqref{eq:ground-track}, καλείται εκ νέου ο πυρήνας υπολογισμού
link-budget (Ενότητες~\ref{sec:meth-terrestrial}-\ref{sec:meth-energy}), και
καταγράφονται όλα τα per-user KPI. Ο βρόχος τερματίζει όταν ο δορυφόρος
γίνει και πάλι αόρατος σε όλους τους χρήστες, αφού έχει προηγουμένως γίνει
ορατός σε τουλάχιστον έναν, καλύπτοντας έτσι ολόκληρη τη διέλευση.

\subsection{Χωρική συσχέτιση shadow fading μεταξύ διαδοχικών βημάτων}
\label{subsec:meth-correlated-sf}

Στο στατικό σενάριο αναφοράς, κάθε κλήση του πυρήνα link-budget αντλεί ένα
νέο, ανεξάρτητο δείγμα LOS/shadow fading ανά ζεύξη BS-χρήστη, κάτι σωστό
όταν κάθε κλήση αντιπροσωπεύει μια εντελώς νέα, ανεξάρτητη πραγματοποίηση
του δικτύου (π.χ. στις επαναλαμβανόμενες εκτελέσεις της
Ενότητας~\ref{sec:results-repeated} ή στα σενάρια του Monte Carlo
generator, Ενότητα~\ref{sec:meth-montecarlo}). Στη χρονική προσομοίωση
όμως οι σταθμοί βάσης και οι χρήστες παραμένουν ακίνητοι από βήμα σε
βήμα, μόνο ο δορυφόρος κινείται, οπότε η επίγεια ζεύξη κάθε χρήστη
βρίσκεται πρακτικά στην ίδια θέση· η επανάληψη ενός τελείως ανεξάρτητου
δείγματος σκίασης κάθε $5$~s δεν αντιστοιχεί σε καμία φυσική μεταβολή του
περιβάλλοντος.

Για τον λόγο αυτό, η χρονική προσομοίωση περνάει την κατάσταση καναλιού
(LOS/NLOS και δείγμα shadow fading ανά ζεύξη) από βήμα σε βήμα, και το νέο
δείγμα shadow fading προκύπτει από το προηγούμενο μέσω εκθετικής
αυτοσυσχέτισης κατά Gudmundson~\cite{gudmundson1991correlation}:
\begin{equation}
\rho(\Delta d) = \exp\!\left(-\frac{\Delta d}{d_{\text{corr}}}\right), \qquad
\text{SF}_k = \rho \cdot \text{SF}_{k-1} + \sqrt{1-\rho^2}\,\sigma_{\text{SF}}\,\mathcal{N}(0,1),
\label{eq:correlated-sf}
\end{equation}
όπου $\Delta d$ η μετατόπιση του χρήστη από το προηγούμενο βήμα ($0$ στην
παρούσα προσομοίωση, αφού οι χρήστες είναι ακίνητοι) και $d_{\text{corr}}$
η "correlation distance in the horizontal plane" του TR~38.901,
Πίνακας~7.5-6 ($37$~m για UMa LOS, $50$~m για UMa NLOS). Η κατάσταση
LOS/NLOS διατηρείται με την ίδια πιθανότητα $\rho$ αντί να
επαναδειγματίζεται σε κάθε βήμα, καθώς το TR~38.901 δεν ορίζει ξεχωριστή
correlation distance για την ίδια την κατηγορική κατάσταση LOS/NLOS·
πρόκειται συνεπώς για μια απλοποίηση συνεπή με τη λογική "drop-based"
παραγωγής παραμέτρων μεγάλης κλίμακας του προτύπου (§7.5, Βήματα 2 και 4).
Με $\Delta d = 0$ (ακίνητος χρήστης), η εξίσωση~\eqref{eq:correlated-sf} δίνει
$\rho = 1$, δηλαδή το shadow fading κάθε επίγειας ζεύξης παραμένει σχεδόν
αναλλοίωτο καθ' όλη τη διέλευση: την αναμενόμενη συμπεριφορά, αφού
τίποτα δεν έχει μετακινηθεί. Η μετάβαση από το ανεξάρτητο στο συσχετισμένο
μοντέλο έχει άμεση, ορατή επίπτωση στη συχνότητα handover που καταγράφεται
(Ενότητα~\ref{sec:results-temporal}).

\section{Monte Carlo generator τυχαιοποιημένων τοπολογιών}
\label{sec:meth-montecarlo}

Για την παραγωγή ενός συνόλου δεδομένων εκπαίδευσης για το μοντέλο
μηχανικής μάθησης (Ενότητα~\ref{sec:meth-ml}), ο πυρήνας link-budget
καλείται επαναληπτικά πάνω σε τυχαιοποιημένες τοπολογίες: αριθμός σταθμών
βάσης ($1$-$4$) και χρηστών ($3$-$10$) ανά σενάριο, θέσεις σταθμών βάσης
γύρω από ένα κέντρο αναφοράς, χρήστες είτε κοντά σε τυχαίο σταθμό βάσης
είτε σε μεγάλη απόσταση (με πιθανότητα $30\%$, ώστε να υπάρχουν και
σενάρια όπου κυριαρχεί ο δορυφόρος), τυχαιοποιημένη γεωμετρία δορυφόρου
(υποδορυφορικό σημείο και υψόμετρο εντός $500$-$600$~km), και τυχαία
επιλογή σεναρίου διάδοσης UMa/UMi. Το ραδιο-configuration (ισχύς, εύρος
ζώνης, μοντέλα ισχύος) παραμένει σταθερό, ίδιο με το στατικό σενάριο
αναφοράς, ώστε το παραγόμενο σύνολο δεδομένων να παραμένει συγκρίσιμο.
Κάθε σειρά του CSV αντιστοιχεί σε έναν χρήστη ενός σεναρίου, και
περιλαμβάνει τόσο τα χαρακτηριστικά/KPI του νικητή κόμβου όσο και
per-candidate διαγνωστικά (καλύτερο BS SNR/distance/path loss, δορυφορικό
SNR/elevation/slant range/path loss), ανεξάρτητα από το ποιος τελικά
κέρδισε, ώστε ένα μοντέλο να μπορεί να μάθει τη \emph{σύγκριση}, όχι
απλώς να διαβάσει το αποτέλεσμα.

\section{Pipeline μηχανικής μάθησης (πρώτη προσέγγιση)}
\label{sec:meth-ml}

Πάνω στο σύνολο δεδομένων της Ενότητας~\ref{sec:meth-montecarlo},
εκπαιδεύονται δύο ταξινομητές -- Λογιστική Παλινδρόμηση και Random Forest
($300$ δέντρα, μέγιστο βάθος $12$) -- με στόχο πρόβλεψης τη μεταβλητή
\texttt{ServingType} (Terrestrial/Satellite), χρησιμοποιώντας
\emph{αποκλειστικά} τα per-candidate χαρακτηριστικά (SNR, path loss,
distance/elevation υποψηφίων), \emph{όχι} τα χαρακτηριστικά του ήδη-
επιλεγμένου νικητή κόμβου, ώστε να αποφευχθεί η τετριμμένη περίπτωση όπου
η ετικέτα διαβάζεται απευθείας από ένα χαρακτηριστικό που ήδη την κωδικοποιεί.
Ο διαχωρισμός εκπαίδευσης/ελέγχου γίνεται ανά \texttt{ScenarioID} (όχι ανά
γραμμή), ώστε χρήστες του ίδιου σεναρίου -- που μοιράζονται γεωμετρία BS/
δορυφόρου -- να μην διαρρέουν πληροφορία μεταξύ συνόλου εκπαίδευσης και
ελέγχου.

Όπως συζητείται αναλυτικά στην Ενότητα~\ref{sec:results-ml}, επειδή η
ετικέτα είναι στην ουσία $\argmax(\text{SNR}_{\text{BS}},
\text{SNR}_{\text{sat}})$ (η ίδια η εξίσωση~\eqref{eq:node-selection}), ένα
μοντέλο στο οποίο δίνονται και τα δύο αυτά SNR αναμένεται να αναπαράγει
τον κανόνα σχεδόν τέλεια. Ο σκοπός αυτής της πρώτης δοκιμής δεν είναι να
αποδείξει προγνωστική ισχύ, αλλά να επαληθεύσει την ορθότητα ολόκληρου του
pipeline δεδομένων (MATLAB $\rightarrow$ CSV $\rightarrow$ Python
$\rightarrow$ εκπαιδευμένο μοντέλο $\rightarrow$ μετρικές), πριν
επενδυθεί προσπάθεια σε πιο απαιτητικούς, μη-τετριμμένους στόχους
πρόβλεψης (βλ.~Κεφάλαιο~\ref{ch:discussion}).
